\begin{center}
    {\LARGE \textbf{2019分析化学}} \\[2ex]
    {\large 时量180分钟 \quad 满分90分}
\end{center}

\vspace{0.5cm}
\noindent\textbf{注意事项:}
\begin{enumerate}[label=\arabic*., parsep=0pt, itemsep=0pt]
    \item 答题前,考生先将自己的姓名、学号、考场号、座位号填写在答题卡上,将条形码准确粘贴在考生信息条形码粘贴区。
    \item 考试结束后,将本试卷与答题纸一并收回。
    \item 回答各个问题时,将答案写在答题纸上,写在本试卷上无效。
\end{enumerate}

\vspace{0.5cm}
\noindent\textbf{一、简答题(42 pts.)}

\begin{enumerate}[label=\arabic*., itemsep=1.5ex]
    \item 采用$\ce{BaSO4}$重量法时，下列各种情况对分析结果会产生什么影响？（偏高、偏低或无影响）
    \begin{center}
    \begin{tabular}{|c|c|c|c|}
    \hline
    \multicolumn{2}{|c|}{有$\ce{Na2SO4}$共沉淀} & \multicolumn{2}{c|}{有$\ce{H2SO4}$共沉淀} \\
    \hline
    测定$\ce{S}$ & 测定$\ce{Ba^{2+}}$ & 测定$\ce{S}$ & 测定$\ce{Ba^{2+}}$ \\
    \hline
     &  &  &  \\
    \hline
    \end{tabular}
    \end{center}
    \item 钢铁中微量铝的测定，可用离子交换法消除铁的干扰。请简要说明实验方法。
    \item 在滴定分析中所用标准溶液浓度不必过大，也不宜过小，其原因是什么。
    \item 写质子条件式：
    \begin{enumerate}[label=(\arabic*), noitemsep, topsep=0pt, partopsep=0pt]
        \item 用$\ce{NaOH}$滴定二氯乙酸($\ce{HA}$,$\mathrm{p}K_a=1.3$)和$\ce{NH4Cl}$混合液中的二氯乙酸至化学计量点时的质子条件式;
        \item 含$0.10\mathrm{mol/L}\ \ce{HCl}$和$0.20\mathrm{mol/L}\ \ce{H2SO4}$的混合溶液的质子条件式。
    \end{enumerate}
    \item $0.001\mathrm{mol/L}\ \ce{Ca(NO3)2}$与$0.001\mathrm{mol/L}\ \ce{(NH4)2C2O4}$等体积混合，在$\mathrm{pH}=1.0$时是否有沉淀产生?
    \\
    $\left[\mathrm{p}K_{\mathrm{sp}}(\ce{CaC2O4})=8.64,\ \mathrm{p}K_{a1}(\ce{H2C2O4})=1.22,\ \mathrm{p}K_{a2}(\ce{H2C2O4})=4.19\right]$
    \item 以双指示剂法进行混合碱($\ce{NaOH}$，$\ce{NaHCO3}$或$\ce{Na2CO3}$)的分析时，取两份体积相同的试液，第一份用酚酞作指示剂，消耗标准$\ce{HCl}$溶液的体积为$V_1$，第二份用甲基橙作指示剂，消耗$\ce{HCl}$溶液的体积为$V_2$，试根据$\ce{HCl}$体积的关系，判断混合碱的组成。
    \begin{center}
    \begin{tabular}{|c|p{6cm}|}
    \hline
    $V_1$、$V_2$关系 & 组成 \\
    \hline
    $V_1=V_2$ &  \\
    \hline
    $V_1=0$ &  \\
    \hline
    $2V_1=V_2$ &  \\
    \hline
    $2V_1>V_2$ &  \\
    \hline
    $2V_1<V_2$ &  \\
    \hline
    \end{tabular}
    \end{center}
    \item 分析天平的称量误差为$\pm0.1\mathrm{mg}$，称样量分别为$0.05\mathrm{g}$、$0.2\mathrm{g}$、$1.0\mathrm{g}$时可能引起的相对误差各为多少?这些结果说明什么问题?
\end{enumerate}

\vspace{0.5cm}
\noindent\textbf{二、计算题(48 pts.)}

\begin{enumerate}[label=\arabic*., itemsep=1.5ex]
    \item (8 pts.) 在$\mathrm{pH}=5.0$的缓冲溶液中，以$2\times10^{-2}\mathrm{mol/L}\ \mathrm{EDTA}$滴定相同浓度的$\ce{Cu^{2+}}$溶液，欲使终点误差在$\pm0.1\%$以内。试通过计算说明可选用$\mathrm{PAN}$指示剂。已知:$\mathrm{pH}=5.0$时$\mathrm{PAN}$的$\mathrm{p}Cu_t=8.8$,$\alpha_{Y(\mathrm{H})}=10^{6.6}$,$\lg K(\ce{CuY})=18.8$。
    \item (10 pts.) 纯明矾中铝的质量分数为$10.77\%$，用某方法测定$9$次结果为($\%$):
    
    $10.74$,$10.77$,$10.77$,$10.77$,$10.81$,$10.82$,$10.73$,$10.86$,$10.81$,求该方法的相对误差。并考察$95\%$置信度的置信区间是否包括真值在内。[$t(0.95,8)=2.31$]
    \item (10 pts.) 用$0.1000\mathrm{mol/L}\ \ce{HNO3}$标准溶液滴定同浓度的$\ce{NH3\cdot H2O}$ ($K_b=1.8\times10^{-5}$)，若指示剂的变色点($\mathrm{p}K_{\mathrm{HIn}}=4.9$)，计算其终点误差。
    \item (10 pts.) 计算下列反应的条件平衡常数(在$1\mathrm{mol/L}\ \ce{HCl}$介质中):
    \\
    $\ce{2Fe^{3+} + 3I^- = 2Fe^{2+} + I3^-}$，当$25\mathrm{mL}\ 0.050\mathrm{mol/L}\ \ce{Fe^{3+}}$与$25\mathrm{mL}\ 0.15\mathrm{mol/L}\ \ce{I^-}$混合后，溶液中残留的$\ce{Fe^{3+}}$还有百分之几?如何才能使$\ce{Fe^{3+}}$定量还原?已知$\varphi^\ominus(\ce{Fe^{3+}/Fe^{2+}})=0.68V$,$\varphi^\ominus(\ce{I3^-/I^-})=0.545V$
    \item (10 pts.) 计算$\ce{ZnS}$在$\mathrm{pH}=8.6$的$0.10\mathrm{mol/L}\ \ce{Na2C2O4}$溶液中的溶解度。[已知$K_{\mathrm{sp}}(\ce{ZnS})=10^{-21.7}$,$\ce{H2C2O4}$的$K_{a1}=10^{-1.2}$,$K_{a2}=10^{-4.2}$;$\ce{Zn(C2O4)3^{4-}}$络合物的$\lg\beta_1\sim\lg\beta_3$分别为$4.9$、$7.6$、$8.2$,$\ce{H2S}$的$K_{a1}=10^{-6.9}$,$K_{a2}=10^{-14.2}$]
\end{enumerate}